%% cover_letter.tex — G7: Real DSO Field-Fault Validation
%% IEEE Transactions on Power Systems
\documentclass[11pt,a4paper]{letter}
\usepackage[T1]{fontenc}
\usepackage[utf8]{inputenc}
\usepackage[margin=2.5cm]{geometry}
\usepackage{microtype}

\signature{Anoop~V.~Eluvathingal\\
           Ph.D.\ Candidate, Dept.\ of Electrical Engineering\\
           Indian Institute of Technology Madras\\
           Chennai 600036, India\\[2pt]
           \texttt{arun.eluvathingal@iitm.ac.in}}

\address{Department of Electrical Engineering\\
         Indian Institute of Technology Madras\\
         Chennai 600036, India}

\begin{document}
\begin{letter}{Editor-in-Chief\\
               \textit{IEEE Transactions on Power Systems}}

\opening{Dear Editor,}

We submit for your consideration the manuscript titled:

\begin{center}
  \textbf{Field Validation of the Self-Adaptive Model-Based Protection
  System on 157 Production IBR Fault Recordings from a 33\,kV DSO
  Network}
\end{center}

\textbf{Motivation.}
Inverter-based resource (IBR) proliferation has profoundly altered
fault-current signatures in distribution networks, yet most protection
algorithms are validated exclusively on synthetic or hardware-in-the-loop
data.  The gap between simulation-domain performance and field performance
is poorly characterised, particularly for high-impedance faults (HIF) and
legacy non-compliant inverters.

\textbf{Contributions.}
This paper makes five distinct contributions:

\begin{enumerate}
  \item \textbf{Field dataset curation (C1).}  We curate and label 157
        disturbance-fault recordings from DSO-A, a real 33\,kV network
        with 12 substations and 47 feeders, covering January 2023 to
        December 2025 (dataset protocol: TR-73).

  \item \textbf{Synthetic-vs-real gap characterisation (C2).}  Systematic
        comparison reveals a $-17.6\%$ HIF detection gap and $+1.1$\,pp
        false-alarm increase when the SAMBPS framework is applied
        unmodified to field data.

  \item \textbf{F1 correction: legacy-IBR $\kappa_n$-trajectory detector
        (C3).}  A sliding-window $\kappa_n$ rise-rate flag identifies
        pre-IEEE-2800 inverter saturation and routes the estimator to a
        one-cycle hold, recovering zone accuracy by $+3.5$\,pp.

  \item \textbf{F2 correction: split-window HIF estimator (C4).}  A
        two-regime negative-sequence estimator handles evolving
        HIF-to-bolted transitions, improving HIF detection by $+7.4$\,pp.

  \item \textbf{F3 correction: shape-factor calibration (C5).}  A scalar
        $\beta$ factor ($\beta_{\rm tri} = 8/\pi^2 \approx 0.811$) corrects
        the asymmetric limiter waveform for unusual inverter profiles,
        closing the remaining detection gap.
\end{enumerate}

Combined, the three corrections (F1--F3) raise zone accuracy from 89.2\%
to 95.5\% and reduce false alarms from 2.1\% to 1.4\%, approaching the
synthetic benchmark of 98.2\% and 1.0\% respectively.

\textbf{Fit to the journal.}
\textit{IEEE Transactions on Power Systems} is the premier venue for
rigorous field-validated protection studies.  This work provides the
first large-scale IBR protection field-validation dataset in a real DSO
context, directly advancing the journal's stated interest in
grid-edge protection and IBR integration.

\textbf{Suggested reviewers.}
\begin{itemize}
  \item Prof.\ Ali Hooshyar (University of Toronto) — IBR protection
  \item Prof.\ Geza Joos (McGill University) — microgrid and distribution protection
  \item Prof.\ Wilsun Xu (University of Alberta) — power quality and HIF detection
  \item Dr.\ Evangelos Farantatos (EPRI) — field validation of protection algorithms
\end{itemize}

This manuscript is not under review elsewhere.  All authors have approved
the submission.  The field dataset was provided under a research agreement
with DSO-A; identifying substation details have been anonymised.

\closing{Sincerely,}
\end{letter}
\end{document}
